\chapter{RDK Bus (RBus) }
\hypertarget{index}{}\label{index}\index{RDK Bus (RBus)@{RDK Bus (RBus)}}
RDK Bus (RBus) is a lightweight, fast and efficient bus messaging system. It allows interprocess communication (IPC) and remote procedure call (RPC) between multiple process running on a hardware device. It supports the creation and use of a data model, which is a hierarchical tree of named objects with properties, events, and methods.

From a developer perspective, there are providers and clients. Providers implement the data model which clients consume.

Providers perform these tasks:
\begin{DoxyItemize}
\item register properties and implement their get and set operations
\item register object and implement their get, set, create, and delete operations
\item register and plublish events
\item register and implement methods which can be called remotely
\end{DoxyItemize}

Consumers perform these tasks:
\begin{DoxyItemize}
\item get and set property values.
\item get, set, create, and delete objects.
\item subscribe and listen to events
\item invoke remote methods
\end{DoxyItemize}

A process can be both a provider and a client. A process can implement many properties, objects, events, and remote methods, and also be a client of the same coming from other providers.

All objects, properties, events, and methods are assigned names by the provider. Each name should be unique across a system, otherwise there will be conflicts in the bus routing.

RBus supports the naming convention defined by TR-\/069, where a name has a hierarchical structure like a file directory structure, with each level of the hierarchy separated by a dot (\textquotesingle{}.\textquotesingle{}), and where object instances and denoted using brace(\{\}).

The RBus API is intended, but not limited, to allow the implementation of a TR-\/181 data model. 